\chapter{Algorytmy przepływowe: Ford--Fulkerson, Edmonds--Karp i Dinic}
\label{ch:tg-flow-algorithms}

W rozdziale \ref{ch:tg-menger} wprowadziliśmy sieci przepływowe,
przepustowości resztowe, przekroje oddzielające $s$ od $t$ i udowodniliśmy twierdzenie
o maksymalnym przepływie i minimalnym przekroju dla
przepustowości całkowitych (twierdzenie \ref{thm:tg-maxflow}).
Obecnie zajmiemy się \emph{algorytmiczną} stroną zagadnienia:
wyborem ścieżek powiększających, ograniczeniami czasu
oraz certyfikatem optymalności.
Rozdział \ref{ch:tg-bfs-dfs} dostarcza BFS i DFS, których
użyjemy jako podstawowych procedur.

\section{Sieci i niezmiennik przepływu}
\label{sec:tg-flow-model}

Niech $D=(V,A)$ będzie skończoną siecią skierowaną bez pętli
z różnymi wierzchołkami $s,t\in V$ oraz
przepustowościami $c:A\to\mathbb R_{\geq0}$.
Oznaczamy $n=|V|$, $m=|A|$.
Przepływ $f$ spełnia $0\leq f(a)\leq c(a)$ dla każdego łuku
oraz prawo zachowania przepływu w $V\setminus\{s,t\}$.
Jego wartość $|f|$ oznacza bilans wypływu ze źródła,
jak w rozdziale \ref{ch:tg-menger}.
Dopuszczamy zarówno łuki $u\to v$, jak i $v\to u$:
są to \emph{dwa różne łuki oryginalne}.
Przechowujemy je oddzielnie, również jeśli ich
przejścia resztowe mają te same końce. Jeśli dopuszczamy
łuki równoległe, każdy łuk ma własny identyfikator,
a zapis $(u,v)$ określa jedynie jego początek i koniec.
Dla $S\subseteq V$ przez pojemność przekroju
$c(S,V\setminus S)$ rozumiemy sumę przepustowości
wszystkich oryginalnych łuków prowadzących z $S$ do
$V\setminus S$. Terminy \emph{pojemność resztowa}
i \emph{przepustowość resztowa} mają tutaj to samo znaczenie.

\begin{definicja}[Przejścia resztowe i ścieżka powiększająca]
\label{def:tg-flow-residual}
Dla oryginalnego łuku $a=(u,v)$ o pojemności $c(a)$
i przepływie $f(a)$ wprowadzamy przejście resztowe
\emph{w przód} $u\to v$ o pojemności $c(a)-f(a)$
oraz przejście \emph{wstecz} $v\to u$ o pojemności $f(a)$.
Zachowujemy informację, z którego oryginalnego łuku
pochodzi dane przejście; zerowe pojemności można pominąć.
\emph{Ścieżką powiększającą} jest prosta skierowana
ścieżka od $s$ do $t$ z dodatnimi pojemnościami
resztowymi na każdym kroku.
Jej \emph{wąskim gardłem} nazywamy minimum tych pojemności.
\end{definicja}

\begin{lemat}[Powiększenie zachowuje dopuszczalność]
\label{lem:tg-flow-augment}
Niech $f$ będzie przepływem dopuszczalnym,
a $P$ ścieżką powiększającą o wąskim gardle
$\varepsilon>0$. Zwiększenie przepływu o
$\varepsilon$ na łukach oryginalnych użytych w przód
i zmniejszenie go o $\varepsilon$ na łukach użytych
wstecz daje przepływ dopuszczalny $f'$,
dla którego $|f'|=|f|+\varepsilon$.
\end{lemat}
\begin{proof}
Każdy krok w przód ma przepustowość resztową
$c(a)-f(a)\geq\varepsilon$, więc po aktualizacji
nadal $f'(a)\leq c(a)$. Na przejściu wstecz użytym przez $P$ mamy
$f(a)\geq\varepsilon$, więc przepływ nie staje się
ujemny. Aktualizacja w przód nie narusza dolnego
ograniczenia, a aktualizacja wstecz nie narusza
górnego ograniczenia przepustowości. Prostota ścieżki gwarantuje, że żaden oryginalny łuk
nie jest aktualizowany dwukrotnie: użycie obu jego
przejść resztowych wymagałoby powtórzenia wierzchołka.
Niezależnie od tego, czy krok resztowy $u\to v$ jest
przejściem w przód czy wstecz, aktualizacja zwiększa
bilans wypływu w $u$ o $\varepsilon$ i zmniejsza go
w $v$ o $\varepsilon$. W każdym wewnętrznym wierzchołku
$P$ zmiany od kroku wchodzącego i wychodzącego
sumują się więc do zera.
W źródle bilans wypływu rośnie o $\varepsilon$,
a w ujściu maleje o tę samą wartość.
\end{proof}

\begin{twierdzenie}[Resztowy certyfikat optymalności]
\label{thm:tg-flow-certificate}
Niech $D=(V,A)$ będzie skończoną siecią skierowaną
bez pętli, $s,t\in V$, $s\ne t$, $c:A\to\mathbb R_{\geq0}$,
a $f$ przepływem dopuszczalnym od $s$ do $t$.
Wówczas następujące warunki są równoważne:
\begin{enumerate}[label=(\roman*)]
\item $f$ ma maksymalną wartość;
\item w sieci resztowej $f$ nie istnieje ścieżka
      powiększająca z $s$ do $t$;
\item istnieje przekrój $S\subset V$ z
      $s\in S$, $t\notin S$, dla którego
      $|f|=c(S,V\setminus S)$.
\end{enumerate}
W takim przypadku $S$ można wybrać jako zbiór
wierzchołków osiągalnych z $s$ w sieci resztowej.
\end{twierdzenie}
\begin{proof}
Jeżeli istnieje ścieżka powiększająca,
lemat \ref{lem:tg-flow-augment} pozwala
ściśle zwiększyć wartość przepływu. Zatem (i)
implikuje (ii). Załóżmy (ii), a przez $S$ oznaczmy
zbiór wierzchołków osiągalnych z $s$
po dodatnich przejściach resztowych.
Wtedy $s\in S$, $t\notin S$.
Każdy oryginalny łuk z $S$ do $V\setminus S$
jest nasycony, ponieważ inaczej istniałoby
dodatnie przejście resztowe na zewnątrz $S$.
Każdy oryginalny łuk w przeciwnym kierunku
ma przepływ zero, ponieważ jego dodatnie
przejście wstecz prowadziłoby poza $S$.
Sumując bilanse wypływu po wierzchołkach $S$,
otrzymujemy (por. lemat \ref{lem:tg-cutbound})
\[
|f|=\sum_{a:S\to V\setminus S}f(a)
-\sum_{a:V\setminus S\to S}f(a).
\]
Łuki wewnętrzne w $S$ wnoszą dwa przeciwne
składniki, a jedynym niezerowym bilansem w $S$
może być bilans źródła. Każdy łuk równoległy
liczymy osobno. Z nasycenia łuków wychodzących
i zerowego przepływu na wchodzących wynika
$|f|=c(S,V\setminus S)$, czyli (iii).
Dla dowolnego innego przepływu dopuszczalnego
$g$ ta sama tożsamość oraz nierówności
$0\leq g(a)\leq c(a)$ dają
$|g|\leq c(S,V\setminus S)=|f|$.
Dowodzi to (iii)$\Rightarrow$(i).
\end{proof}

\begin{wniosek}[Twierdzenie o maksymalnym przepływie
i minimalnym przekroju dla przepustowości rzeczywistych]
\label{cor:tg-flow-real-mincut}
Niech $D=(V,A)$ będzie skończoną siecią skierowaną
bez pętli, $s,t\in V$, $s\ne t$ i
$c:A\to\mathbb R_{\geq0}$. Wówczas istnieje
przepływ maksymalny od $s$ do $t$ i zachodzi
\[
\max_f |f|=\min_{\substack{s\in S\subseteq V\\t\notin S}}
c(S,V\setminus S).
\]
Maksimum bierzemy po wszystkich przepływach
dopuszczalnych od $s$ do $t$.
\end{wniosek}
\begin{proof}
Zbiór wszystkich dopuszczalnych przepływów
jest niepusty, ponieważ zawiera przepływ zerowy.
Jest przecięciem prostopadłościanu
$\prod_{a\in A}[0,c(a)]$ z podprzestrzenią
liniową wyznaczoną przez równania zachowania przepływu.
Jest zatem domkniętym i ograniczonym podzbiorem
skończenie wymiarowej przestrzeni $\mathbb R^A$,
a więc jest zwarty.
Liniowa funkcja $f\mapsto |f|$ osiąga na nim
maksimum. Dla przepływu maksymalnego
twierdzenie \ref{thm:tg-flow-certificate}
dostarcza przekroju o równej pojemności,
a lemat \ref{lem:tg-cutbound} daje nierówność
dla wszystkich innych przepływów i przekrojów.
\end{proof}

\begin{twierdzenie}[Całkowity rozkład przepływu i drogi rozłączne krawędziowo]
\label{thm:tg-flow-integral-paths}
Niech $D=(V,A)$ będzie skończoną siecią skierowaną bez pętli,
z różnymi wierzchołkami $s,t$ i nieujemnymi całkowitymi
przepustowościami $c(a)$. Niech $f$ będzie dopuszczalnym
przepływem całkowitoliczbowym o nieujemnej wartości
$k=|f|\geq0$.
Wówczas istnieją skierowane proste ścieżki
$P_1,\ldots,P_\ell$ od $s$ do $t$,
dodatnie liczby całkowite $q_1,\ldots,q_\ell$
oraz skończona rodzina skierowanych cykli prostych
$C_1,\ldots,C_h$ o dodatnich całkowitych
współczynnikach $b_j$, takie że
\[
 f=\sum_{i=1}^{\ell}q_i\chi_{P_i}
   +\sum_{j=1}^{h}b_j\chi_{C_j},
 \qquad \sum_{i=1}^{\ell}q_i=k.
\]
Dla łuku $a\in A$ symbole $\chi_P(a)$ i $\chi_C(a)$
oznaczają funkcje charakterystyczne odpowiednich
zbiorów łuków. Jeśli ponadto wszystkie przepustowości
wynoszą $0$ albo $1$, niech $D_+$ oznacza podgraf
o zbiorze wierzchołków $V$, utworzony z łuków
o przepustowości $1$. Wówczas istnieje rodzina
$k$ parami rozłącznych krawędziowo ścieżek
skierowanych od $s$ do $t$ w $D_+$.
Maksymalna liczba takich ścieżek w $D_+$
jest równa najmniejszej liczbie łuków $D_+$,
których usunięcie niszczy wszystkie drogi skierowane
od $s$ do $t$ w $D_+$.
\end{twierdzenie}
\begin{proof}
Przypuśćmy najpierw, że $k>0$, i rozważmy
podgraf złożony z łuków, na których $f(a)>0$.
Pokażemy, że zawiera on ścieżkę skierowaną
od $s$ do $t$. W przeciwnym razie zbiór $S$
wierzchołków osiągalnych z $s$ po takich łukach
nie zawierałby $t$, a żaden łuk
$S\to V\setminus S$ nie niósłby dodatniego
przepływu. Z tożsamości przepływu przez przekrój
otrzymalibyśmy
\[
k=\sum_{a\colon S\to V\setminus S}f(a)
  -\sum_{a\colon V\setminus S\to S}f(a)\leq0,
\]
sprzeczność. Wybierzmy więc prostą ścieżkę $P$
z dodatnim przepływem na każdym łuku. Oznaczmy
bieżącą wartość przepływu przez $r>0$ i połóżmy
\[
q=\min\{r,\min_{a\in P}f(a)\}.
\]
Liczba $q$ jest całkowita i dodatnia.
Warunek $q\leq r$ zapobiega zmianie znaku wartości
przepływu przy usuwaniu ścieżki, nawet gdy
w sieci występuje cyrkulacja.

Odjęcie $q\chi_P$ zachowuje
nieujemność i ograniczenia przepustowości:
na łukach $P$ mamy $q\leq f(a)$, a na pozostałych
przepływ pozostaje bez zmian. W każdym wewnętrznym
wierzchołku $P$ zmniejszamy wpływ i wypływ
o tę samą wartość. W źródle odejmujemy $q$
od bilansu wypływu, w ujściu zaś dodajemy $q$.
Otrzymany przepływ jest dopuszczalny i ma
wartość $r-q$.
Powtarzamy ten krok, aż wartość stanie się zerowa.
Każdy krok albo kończy tę fazę (gdy $q=r$),
albo zeruje przepływ na co najmniej jednym łuku.
Wartości na łukach nie rosną, więc liczba kroków
jest skończona.

Pozostały nieujemny przepływ całkowitoliczbowy
o wartości zero jest cyrkulacją. Bilans wypływu
w $s$ jest zerowy z definicji wartości, a suma
bilansów wszystkich wierzchołków wynosi zero,
gdyż każdy łuk wnosi dwa przeciwne składniki.
Ponieważ pozostałe wierzchołki poza $t$ mają
bilans zero, także $t$ ma bilans zero.
Jeżeli cyrkulacja jest niezerowa, wybierzmy
łuk o dodatnim przepływie i rozpocznijmy nim spacer.
W każdym osiągniętym wierzchołku dodatni wpływ
wymusza dodatni wypływ. Spacer można więc kontynuować
aż do pierwszego powtórzenia wierzchołka; fragment
między jego dwoma wystąpieniami jest cyklem prostym.
Odejmujemy minimalny dodatni przepływ na łukach
tego cyklu. W każdym jego wierzchołku wpływ
i wypływ zmniejszają się o tę samą liczbę,
więc pozostały przepływ nadal jest cyrkulacją. Także tutaj każdy krok
zeruje co najmniej jeden łuk, więc
po skończonej liczbie kroków otrzymujemy zero.
Dodając odjęte składniki, odtwarzamy początkowy
przepływ $f$ na każdym łuku. Suma współczynników
ścieżek wynosi $k$, ponieważ ich odjęcie zmniejszyło
wartość z $k$ do zera, a cykle mają bilans zerowy.
Dla $k=0$ rodzina ścieżek jest pusta.

Gdy wszystkie $c(a)\in\{0,1\}$,
na każdym łuku $f(a)\leq1$. Stąd
wszystkie współczynniki $q_i,b_j$
są równe $1$, a ścieżki $P_i$
nie mają wspólnych łuków.
Ponieważ żadna dodatnia część przepływu
nie przechodzi przez łuk o przepustowości $0$,
wszystkie te ścieżki leżą w $D_+$.
Ich liczba wynosi $k$.
Odwrotnie, każda rodzina $k$
parami rozłącznych krawędziowo
ścieżek w $D_+$ wyznacza dopuszczalny
przepływ o wartości $k$.
Każda taka ścieżka musi przejść z $S$ do
$V\setminus S$ co najmniej jednym łukiem $D_+$.
Z rozłączności krawędziowej wynika, że wybrane
łuki dla różnych ścieżek są różne. Liczba
ścieżek nie przekracza więc pojemności przekroju.
Dla przepustowości $0$ lub $1$
pojemność przekroju jest dokładnie
liczbą jego łuków należących do $D_+$.
Twierdzenie o maksymalnym przepływie
i minimalnym przekroju
\ref{cor:tg-flow-real-mincut}
daje zatem przekrój o liczbie takich łuków
równej maksymalnej wartości
przepływu; maksimum osiągamy
przepływem całkowitoliczbowym
na mocy algorytmu Forda--Fulkersona
z Twierdzenia~\ref{thm:tg-flow-ff-integer}.
Po usunięciu tych łuków
nie ma drogi od $s$ do $t$ w $D_+$.
Z drugiej strony każdy zbiór usuwanych łuków,
który niszczy wszystkie takie drogi, musi zawierać
co najmniej jeden łuk z każdej ścieżki rodziny
otrzymanej z maksymalnego przepływu.
Rozłączność tej rodziny wymusza co najmniej tyle
usuwanych łuków, ile wynosi wartość maksimum.
Dowodzi to równości z minimalną licznością
dowolnego zbioru rozcinającego.
Nie można w tym wniosku liczyć dróg
używających łuków o przepustowości $0$.
\end{proof}

\begin{przyklad}[Cyrkulacja przechodząca przez źródło i ujście]
Niech sieć ma dwa przeciwne łuki $s\to t$ oraz $t\to s$
o przepustowościach odpowiednio $2$ i $1$.
Dla przepływu nasycającego oba łuki wartość wynosi $1$.
Ścieżka $s\to t$ niesie jednak dwie jednostki.
Odjęcie dwóch jednostek dałoby wartość $-1$.
Odejmujemy więc jedną jednostkę, a pozostały przepływ
jest cyrkulacją o wartości zero na cyklu $s\to t\to s$.
\end{przyklad}

\begin{przyklad}[Dwie rozłączne drogi mimo łuku poprzecznego]
\label{ex:tg-flow-integral-paths}
Niech sieć ma łuki
$s\to a$, $s\to b$, $a\to t$, $b\to t$
i $a\to b$, wszystkie o przepustowości $1$.
Istnieją dwie rozłączne krawędziowo drogi
$s\to a\to t$ i $s\to b\to t$.
Przekrój $S=\{s\}$ ma dokładnie dwa łuki,
więc maksymalna liczba rozłącznych dróg
jest równa $2$. W szczególności
łuk poprzeczny $a\to b$ nie zmienia
tej liczby.
\end{przyklad}

\section{Ford--Fulkerson: wybór dowolnej ścieżki}
\label{sec:tg-flow-ff}

Schemat Forda--Fulkersona rozpoczyna od
przepływu zerowego. Dopóki istnieje ścieżka
powiększająca, wybiera \emph{dowolną} taką
ścieżkę i przesyła przez nią jej wąskie gardło.
Można szukać ścieżki za pomocą DFS, BFS
albo inną regułą. Po zakończeniu optymalność
wynika z twierdzenia \ref{thm:tg-flow-certificate}.

\begin{twierdzenie}[Zakończenie przy przepustowościach całkowitych]
\label{thm:tg-flow-ff-integer}
Jeżeli wszystkie przepustowości są nieujemnymi
liczbami całkowitymi, to schemat Forda--Fulkersona
z dowolnym wyborem ścieżki powiększającej,
rozpoczęty od przepływu zerowego,
zatrzymuje się z przepływem maksymalnym.
Jego liczba powiększeń jest nie większa niż
$C_s:=\sum_{a:\,s\to v}c(a)$, gdzie sumujemy
po wszystkich łukach wychodzących z $s$.
Przy znalezieniu ścieżki za pomocą DFS
czas wynosi $O((n+m)(C_s+1))$.
Jest to ograniczenie pseudowielomianowe:
zależy od wartości liczbowej $C_s$, a nie od
długości jej zapisu binarnego.
\end{twierdzenie}
\begin{proof}
Z lematu \ref{lem:tg-flow-augment} przepływy
pozostają dopuszczalne. Indukcja pokazuje, że
są całkowitoliczbowe, ponieważ wszystkie
pojemności resztowe i wąskie gardła
są całkowite. Każde powiększenie zwiększa
wartość o co najmniej $1$, a ta jest
ograniczona przez $C_s$, na mocy
lematu \ref{lem:tg-cutbound} dla $S=\{s\}$.
Po najwyżej $C_s$ powiększeniach algorytm
musi się zatrzymać. Wówczas brak ścieżki
resztowej daje optymalność.
Budowa list resztowych i jedno ich przeszukanie
kosztują $O(n+m)$, a aktualizacja prostej ścieżki
kosztuje $O(n)$. Do co najwyżej $C_s$ udanych
przeszukań dochodzi ostatnie, nieudane przeszukanie
potwierdzające brak ścieżki. Stąd zapowiedziane
ograniczenie liczby operacji arytmetycznych.
\end{proof}

\begin{przyklad}[Wiele powiększeń przy dużych przepustowościach]
\label{ex:tg-flow-ff-pseudopoly}
Niech $N\geq1$ będzie liczbą całkowitą. W sieci
z wierzchołkami $s,a,b,t$ przyjmijmy przepustowość
$N$ na czterech łukach $s\to a$, $s\to b$,
$a\to t$, $b\to t$ oraz przepustowość $1$
na łuku $a\to b$. Innych łuków nie ma.

Wybierajmy naprzemiennie ścieżki $s\to a\to b\to t$
oraz $s\to b\to a\to t$, w drugiej z nich używając
przejścia resztowego wstecz względem $a\to b$.
Po $2j$ powiększeniach ($0\leq j\leq N$) na każdym
z czterech zewnętrznych łuków płynie $j$ jednostek,
a na $a\to b$ przepływ jest zerowy. Istotnie, dla
$j<N$ pierwsza ścieżka przesyła jedną jednostkę,
nasycając $a\to b$. Druga przesyła także jedną
jednostkę i zeruje przepływ na tym łuku.
Przepływ na wszystkich zewnętrznych łukach zwiększa
się w rezultacie z $j$ do $j+1$.

Otrzymujemy dokładnie $2N$ powiększeń i maksymalny
przepływ wartości $2N$, równej pojemności przekroju
$\{s\}$. Przy $N=2^k$ długość zapisu przepustowości
wynosi $O(k+1)$, lecz liczba powiększeń jest
wykładnicza względem $k$. Natomiast wybór dwóch
ścieżek $s\to a\to t$ i $s\to b\to t$ pozwala
osiągnąć ten sam wynik po dwóch powiększeniach.
Przykład pokazuje znaczenie reguły wyboru ścieżki.
\end{przyklad}

\begin{uwaga}[Pojemności wymierne i niewymierne]
Jeśli wszystkie przepustowości są wymierne,
wybieramy dodatnią liczbę całkowitą $L$ będącą
wspólną wielokrotnością ich mianowników.
Zastąpienie $c$ przez $Lc$ i $f$ przez $Lf$
mnoży wszystkie przepustowości resztowe oraz
wąskie gardła przez $L$, nie zmieniając zbioru
dodatnich przejść. Każdemu przebiegowi w sieci
wyjściowej odpowiada zatem przebieg całkowitoliczbowy.
Otrzymujemy zakończenie po co najwyżej $LC_s$
powiększeniach, lecz nie jest to ograniczenie
silnie wielomianowe.
Przy rzeczywistych pojemnościach niewymiernych
niektóre sposoby wybierania ścieżek
mogą prowadzić do nieskończonego ciągu powiększeń.
Nie można zatem utożsamiać dowolnego wyboru ścieżek
z bezwarunkowo kończącym się algorytmem.

\end{uwaga}

\section{Po co potrzebne są przejścia wsteczne?}
\label{sec:tg-flow-backward}

\begin{przyklad}[Zmiana wcześniejszej decyzji]
\label{ex:tg-flow-cancel}
Rozważmy sieć o wierzchołkach $\{s,a,b,t\}$
i pojemnościach $1$ na łukach
$s\to a$, $s\to b$, $a\to b$, $a\to t$ i $b\to t$.
Jeśli pierwszy wybór to ścieżka
$s\to a\to b\to t$, wysyłamy jednostkę.
Pozornie nie ma już wolnej drogi z $s$ do $t$
po samych oryginalnych łukach.
W sieci resztowej istnieje jednak ścieżka
$s\to b\to a\to t$,
gdzie krok $b\to a$ jest przejściem
\emph{wstecz} odpowiadającym oryginalnemu
łukowi $a\to b$. Drugie powiększenie zeruje
przepływ $a\to b$ i daje ostatecznie
jednostkę na $s\to a\to t$
oraz jednostkę na $s\to b\to t$.
Wartość przepływu wynosi $2$ i jest maksymalna,
ponieważ przekrój $\{s\}$ ma pojemność $2$.
\end{przyklad}

\begin{figure}[htbp]
\centering
\begin{tikzpicture}[>=Stealth,scale=.95]
\node[tgBlueVertex] (s) at (-3,0) {$s$};
\node[tgTealVertex] (a) at (0,1.35) {$a$};
\node[tgTealVertex] (b) at (0,-1.35) {$b$};
\node[tgBlueVertex] (t) at (3,0) {$t$};
\draw[tgTree,->] (s)--node[above left,font=\small] {$1$} (a);
\draw[tgTree,->] (s)--node[below left,font=\small] {$1$} (b);
\draw[tgTree,->] (a)--node[above right,font=\small] {$1$} (t);
\draw[tgTree,->] (b)--node[below right,font=\small] {$1$} (t);
\draw[tgAccent,->] (a)--node[right,font=\small] {$1$} (b);
\end{tikzpicture}
\caption{Sieć o przepływie maksymalnym $2$.
Łuk $a\to b$ może mieć dodatni przepływ na pierwszym
etapie, lecz w rozwiązaniu maksymalnym jest zerowy.
Przejście wsteczne $b\to a$ umożliwia korektę.}
\label{fig:tg-flow-backward}
\end{figure}

\section{Edmonds--Karp: najkrótsza ścieżka powiększająca}
\label{sec:tg-flow-ek}

Algorytm Edmondsa--Karpa to Ford--Fulkerson,
w którym każdą ścieżkę powiększającą wyznaczamy
przez BFS w aktualnej sieci resztowej.
Wszystkie przejścia resztowe traktujemy
jako łuki długości $1$, niezależnie
od ich przepustowości. Pamiętamy
identyfikator \emph{konkretnego} przejścia,
aby rozróżnić łuki przeciwne i równoległe.

\begin{lemat}[Monotoniczność odległości resztowych]
\label{lem:tg-flow-ek-levels}
Niech $d_f(v)$ oznacza minimalną liczbę łuków
na ścieżce z $s$ do $v$ w sieci resztowej
przepływu $f$, z wartością $+\infty$ dla
wierzchołków nieosiągalnych.
Jeśli $f'$ powstał przez powiększenie
wzdłuż ścieżki najkrótszej w liczbie łuków,
to $d_{f'}(v)\geq d_f(v)$ dla każdego $v\in V$.
\end{lemat}
\begin{proof}
Przypuśćmy, że teza jest fałszywa,
i wybierzmy wierzchołek $v$ o najmniejszej
nowej odległości $d_{f'}(v)$ spośród tych,
których odległość się zmniejszyła.
Wybierzmy poprzednik $u$ na najkrótszej
nowej ścieżce do $v$. Wtedy
$d_{f'}(u)=d_{f'}(v)-1$ oraz
$d_{f'}(u)\geq d_f(u)$ na mocy
minimalności wyboru $v$.
Wierzchołek $v$ nie jest źródłem, a $d_f(u)$
jest skończone, ponieważ $d_f(u)\leq d_{f'}(u)<+\infty$.
Jeśli konkretne przejście $u\to v$ miało dodatnią
przepustowość już przed powiększeniem,
mamy $d_f(v)\leq d_f(u)+1
\leq d_{f'}(u)+1=d_{f'}(v)$,
sprzeczność (w tym gdy $d_f(v)=+\infty$).
Jeśli przejście pojawiło się dopiero
po powiększeniu, jest sprzężonym przejściem
odwrotnym (w przód lub wstecz)
do pewnego przejścia $v\to u$
wykorzystanego na wybranej najkrótszej
ścieżce powiększającej. Stąd
$d_f(u)=d_f(v)+1$, a zatem
$d_{f'}(v)=d_{f'}(u)+1
\geq d_f(u)+1=d_f(v)+2$,
znów sprzeczność.
\end{proof}

\begin{twierdzenie}[Edmonds--Karp]
\label{thm:tg-flow-ek}
Dla każdej skończonej sieci o rzeczywistych
nieujemnych przepustowościach algorytm
Edmondsa--Karpa kończy działanie
po $O(nm)$ powiększeniach
i zwraca przepływ maksymalny
wraz z przekrojem minimalnym.
Dla reprezentacji listowej czas wynosi
$O\bigl(nm(n+m)+n+m\bigr)$;
w szczególności $O(nm^2)$,
gdy $m\geq n$.
Nie wymaga całkowitości przepustowości.
\end{twierdzenie}
\begin{proof}
Każde powiększenie nasyca przynajmniej
jedno przejście resztowe $e=(u,v)$.
Na ścieżce wyznaczonej przez BFS
zachodzi $d_f(v)=d_f(u)+1$.
Aby \emph{to samo przejście} $e$
mogło być nasycone powtórnie,
jego pojemność musi zostać
wcześniej zwiększona przez
powiększenie wykorzystujące
przejście odwrotne $v\to u$.
Niech $d_1$ oznacza odległości tuż przed pierwszym
nasyceniem $e$, a $d_2$ odległości tuż przed
późniejszym użyciem jego sprzężonego przejścia.
Mamy $d_1(v)=d_1(u)+1$ oraz $d_2(u)=d_2(v)+1$.
Z lematu \ref{lem:tg-flow-ek-levels} otrzymujemy
$d_2(v)\geq d_1(v)$, a stąd
$d_2(u)\geq d_1(u)+2$.
Przed następnym nasyceniem $e$ odległość do $u$
nie może być mniejsza od $d_2(u)$.
Zatem pomiędzy kolejnymi nasyceniami tego samego
rekordu resztowego wzrasta o co najmniej $2$.
Dopóki $u$ jest osiągalny,
odległość ta jest mniejsza niż $n$.
Dlatego to samo przejście
może stać się krytyczne
co najwyżej $O(n)$ razy.
Ponieważ przejść resztowych
z przypisanym oryginalnym łukiem
jest co najwyżej $2m$,
liczba powiększeń wynosi $O(nm)$.
Każde BFS kosztuje $O(n+m)$, a odtworzenie
ścieżki z identyfikatorów przejść i aktualizacja
przepływu kosztują $O(n)$. Uwzględniając ostatnie
BFS bez ścieżki, otrzymujemy
$O((nm+1)(n+m))$ operacji.
Dla $m\geq n$ mamy $n+m\leq 2m$,
co daje $O(nm^2)$.
Algorytm kończy się więc
także dla przepustowości rzeczywistych,
a twierdzenie \ref{thm:tg-flow-certificate}
daje wartość maksymalną
i przekrój minimalny.
\end{proof}

\section{Algorytm Dinica: sieć warstwowa i przepływ blokujący}
\label{sec:tg-flow-dinic}

Algorytm Dinica ogranicza liczbę uruchomień BFS
przez wykonywanie wielu powiększeń przepływu
w obrębie jednej sieci warstwowej.

\begin{definicja}[Sieć warstwowa]
\label{def:tg-flow-level}
Dla aktualnego przepływu $f$
wyznaczamy BFS od $s$ w sieci resztowej;
niech $\ell(v)$ oznacza liczbę łuków
najkrótszej resztowej ścieżki $s\to v$,
a dla wierzchołków nieosiągalnych niech
$\ell(v)=+\infty$.
\emph{Sieć warstwowa} zawiera
tylko dodatnie przejścia resztowe
$u\to v$ spełniające
$\ell(u)<+\infty$ oraz $\ell(v)=\ell(u)+1$.
Warunek skończoności wyklucza łuki między
nieosiągalnymi wierzchołkami.
Gdy $t$ jest nieosiągalny, algorytm kończy działanie.
\end{definicja}

\begin{definicja}[Przepływ blokujący]
\label{def:tg-flow-blocking}
Niech $f$ będzie przepływem dopuszczalnym,
a $L$ jego siecią warstwową z przepustowościami
resztowymi ustalonymi na początku fazy.
Przepływ dopuszczalny $g$ od $s$ do $t$ w $L$
nazywamy \emph{blokującym}, jeśli każda ścieżka
od $s$ do $t$ w $L$ zawiera przejście nasycone
przez $g$. Równoważnie: po zastosowaniu $g$
do przepływu $f$ nie pozostaje żadna ścieżka
od $s$ do $t$ złożona z dodatnich przejść
pierwotnej sieci warstwowej. Nie wymagamy
maksymalności $g$ w sieci warstwowej.
\end{definicja}

Jedna \emph{faza Dinica} tworzy sieć warstwową
przez BFS, a następnie wyznacza przepływ
blokujący, uruchamiając DFS po krawędziach
prowadzących o jedną warstwę wyżej.
Dla każdego wierzchołka przechowujemy
wskaźnik następnego niezbadanego przejścia
(\emph{current arc}). Po nasyceniu przejścia
lub wykazaniu, że nie daje już drogi do $t$,
wskaźnik przesuwamy dalej.
W kolejnej fazie wyznaczamy nowe warstwy.

\begin{lemat}[Wzrost odległości do ujścia]
\label{lem:tg-flow-dinic-distance}
Po zakończeniu fazy Dinica, jeśli
$t$ pozostaje osiągalny resztowo,
to jego nowa odległość od $s$
jest \emph{ściśle większa}
od odległości na początku fazy.
\end{lemat}
\begin{proof}
Niech $\ell$ oznacza odległości ustalone na początku
fazy. Każde początkowe dodatnie przejście $u\to v$
o skończonej $\ell(u)$ spełnia
$\ell(v)\leq\ell(u)+1<+\infty$ na mocy własności BFS.
Nowe dodatnie przejścia mogą powstać jedynie jako
sprzężone odwrotności użytych przejść warstwowych,
a zatem prowadzą z warstwy $j+1$ do warstwy $j$.
Własność $\ell(v)\leq\ell(u)+1$ dla skończonej
$\ell(u)$ zachodzi więc przez całą fazę.

Wzdłuż dowolnej resztowej ścieżki po fazie,
zaczynającej się w $s$, wszystkie stare odległości
są skończone (indukcja od $\ell(s)=0$),
a każdy krok zwiększa starą odległość o co najwyżej
jeden. Każda taka ścieżka do $t$ ma zatem długość
co najmniej $\ell(t)$. Jeśli ma dokładnie $\ell(t)$
przejść, to każde z nich musi zwiększać starą
odległość dokładnie o jeden. Żadne nie jest nowo
utworzonym przejściem odwrotnym.
Byłaby to zatem dodatnia
ścieżka w pozostałej sieci warstwowej,
sprzecznie z definicją
przepływu blokującego.
\end{proof}

\begin{twierdzenie}[Zakończenie i złożoność Dinica]
\label{thm:tg-flow-dinic}
Dla skończonej sieci o $n\geq2$ wierzchołkach,
$m$ łukach i nieujemnych przepustowościach
rzeczywistych algorytm Dinica
kończy działanie po najwyżej $n-1$
fazach z ujściem osiągalnym
i zwraca przepływ maksymalny.
W standardowej implementacji
z listami sąsiedztwa i wskaźnikami
\emph{current arc} czas wynosi
$O\bigl(n^2(m+1)\bigr)$,
a pamięć $O(n+m)$.
Dla $m\geq1$ odpowiada to
klasycznemu oszacowaniu $O(n^2m)$.
\end{twierdzenie}
\begin{proof}
Odległość $\ell(t)$ do
osiągalnego ujścia należy
do $\{1,\ldots,n-1\}$.
Na mocy lematu
\ref{lem:tg-flow-dinic-distance}
rośnie ściśle po każdej fazie,
zatem liczba faz z osiągalnym
ujściem nie przekracza $n-1$.
Po ich zakończeniu brak ścieżki
powiększającej daje optymalność.

Uzasadnijmy koszt jednej fazy.
Budowa warstw przez BFS
zajmuje $O(n+m)$.
W obrębie sieci warstwowej dodatnie pojemności
przejść warstwowych mogą jedynie maleć.
Jeżeli DFS wyczerpie wszystkie drogi z danego
wierzchołka do $t$, późniejsze powiększenia nie
mogą ich odtworzyć. Wskaźnik można zatem trwale
przesunąć po nasyceniu przejścia lub po nieudanym
wywołaniu dla jego końca. Łączna liczba przesunięć
wskaźników i nieudanych prób jest $O(n+m)$.
Udane zejścia DFS liczymy osobno, w koszcie ścieżek
powiększających. Ponawiamy DFS od $s$, aż nie znajdzie
żadnej dodatniej ścieżki warstwowej; wówczas z definicji
uzyskaliśmy przepływ blokujący.
Każde skuteczne powiększenie
w obrębie danej fazy
nasyca co najmniej jedno
przejście warstwowe,
które nie odzyskuje
pojemności w tej fazie:
przejścia odwrotne
prowadzą do wcześniejszej
warstwy i nie są używane.
Dla każdego oryginalnego łuku co najwyżej jedno
z jego dwóch sprzężonych przejść może być warstwowe,
gdyż drugie prowadzi w przeciwnym kierunku.
Jest więc co najwyżej $m$ takich powiększeń,
z których każde wymaga
$O(n)$ aktualizacji
na ścieżce. Otrzymujemy
$O(nm+n+m)$ czasu
na fazę oraz
$O(n+m)$ pamięci.
Pomnożenie przez co najwyżej $n-1$ faz i dodanie
kosztu końcowego BFS daje
$O(n(nm+n+m)+n+m)=O(n^2(m+1))$.
Tablice poziomów i wskaźników oraz stos DFS
zajmują $O(n)$ pamięci, a listy sprzężonych
przejść $O(n+m)$.
\end{proof}

\section{Certyfikat, przykład liczbowy i uwagi implementacyjne}
\label{sec:tg-flow-example}

\begin{przyklad}[Przepływ $5$ i przekrój minimalny]
\label{ex:tg-flow-cut-five}
Niech $V=\{s,a,b,t\}$
oraz przepustowości
$c(s,a)=3$, $c(s,b)=2$,
$c(a,b)=1$, $c(a,t)=2$,
$c(b,t)=3$.
Przepływ
$f(s,a)=3$, $f(s,b)=2$,
$f(a,b)=1$, $f(a,t)=2$,
$f(b,t)=3$
jest dopuszczalny:
w $a$ wpływa $3$,
a wypływa $1+2=3$;
w $b$ wpływa $2+1=3$,
a wypływa $3$.
Ma wartość $5$.
Przekrój $S=\{s\}$ ma pojemność
$c(s,a)+c(s,b)=5$.
Wobec lematu \ref{lem:tg-cutbound}
przepływ jest maksymalny,
a ten przekrój minimalny.
\end{przyklad}

\begin{figure}[htbp]
\centering
\begin{tikzpicture}[>=Stealth,scale=.92]
\node[tgBlueVertex] (s) at (-3,0) {$s$};
\node[tgTealVertex] (a) at (0,1.45) {$a$};
\node[tgTealVertex] (b) at (0,-1.45) {$b$};
\node[tgBlueVertex] (t) at (3,0) {$t$};
\draw[tgTree,->] (s)--node[above left,font=\small] {$3$} (a);
\draw[tgTree,->] (s)--node[below left,font=\small] {$2$} (b);
\draw[tgTree,->] (a)--node[above right,font=\small] {$2$} (t);
\draw[tgTree,->] (b)--node[below right,font=\small] {$3$} (t);
\draw[tgEdge,->] (a)--node[right,font=\small] {$1$} (b);
\draw[tgAccent,densely dashed] (-2.48,-1.8)--(-2.48,1.8);
\end{tikzpicture}
\caption{Sieć o maksymalnym przepływie $5$.
Liczby przy łukach są ich przepustowościami.
Przerywana pionowa linia oddziela $S=\{s\}$;
przecinają ją łuki o łącznej pojemności $5$.}
\label{fig:tg-flow-cut-five}
\end{figure}

\begin{stwierdzenie}[Sposób odzyskania przekroju minimalnego]
\label{prop:tg-flow-recover-cut}
Po zakończeniu Edmondsa--Karpa lub Dinica
wykonajmy BFS z $s$ po dodatnich przejściach
resztowych i oznaczmy zbiór osiągalnych
wierzchołków przez $S$.
Wówczas $S$ wyznacza przekrój minimalny,
a jego pojemność równa jest
wartości znalezionego przepływu.
Koszt tej czynności jest $O(n+m)$.
\end{stwierdzenie}
\begin{proof}
Brak resztowej ścieżki do $t$
i twierdzenie \ref{thm:tg-flow-certificate}
dają minimalność i równość.
BFS przegląda co najwyżej
$2m$ przejść i $n$ wierzchołków.
\end{proof}

\begin{uwaga}[Reprezentacja i arytmetyka]
W implementacji dodajemy dla każdego łuku
dwa sprzężone rekordy resztowe:
rekord w przód i rekord wstecz,
z indeksami wskazującymi wzajemnie
na siebie. Zapobiega to błędom
przy równoległych przejściach i
oryginalnych łukach przeciwnych.
Pojemności całkowite warto przechowywać
w typie liczbowym o odpowiednim zakresie;
dla dokładnych ułamków można używać
liczb wymiernych. Przy arytmetyce
zmiennopozycyjnej porównywanie
pojemności z zerem i błędy zaokrągleń
wymagają odrębnej analizy numerycznej.
Wszystkie dowody powyżej odnoszą się
do \emph{dokładnej} arytmetyki rzeczywistej.
\end{uwaga}

\section{Porównanie algorytmów}
\label{sec:tg-flow-summary}

Dla $n=|V|$ i $m=|A|$
otrzymaliśmy następujące rozróżnienie:
\begin{itemize}[leftmargin=2em,itemsep=3pt]
\item \textbf{Ford--Fulkerson}: prosta idea;
      dla przepustowości całkowitych
      liczba powiększeń jest ograniczona przez wartość
      maksymalnego przepływu, a dla dowolnych
      pojemności rzeczywistych bez
      ograniczenia wyboru ścieżek
      może nie kończyć działania.
\item \textbf{Edmonds--Karp}: wybór ścieżki BFS,
      gwarantowana wielomianowa liczba
      powiększeń dla dowolnych pojemności
      rzeczywistych; czas
      $O(nm(n+m)+n+m)$.
\item \textbf{Dinic}: BFS warstwowy i
      przepływ blokujący;
      czas $O(n^2(m+1))$
      i pamięć $O(n+m)$.
\end{itemize}

Podstawą poprawności każdego kończącego
się algorytmu jest ten sam certyfikat:
brak ścieżki powiększającej
i przekrój $S$ o pojemności
równej wartości przepływu.
Przy wyborze implementacji
istotne są także liczby łuków,
przepustowości i struktura sieci.
